\section{Research target, contribution, and conventions}
\label{asj:sec:scope}

\subsection{A specific repository question}
The manuscript's Hurwitz-jet chapters organize parameter derivatives, spectral
derivatives, generalized Stieltjes constants, and negative-order polygamma in one
calculus \cite{repo}. A later report already proves a normalized-exponent polynomial
law for polynomial spectral multiplicities, including a polynomial times a power of
a logarithm \cite{prior}. Its question Q10 asks for arithmetic and nonpolynomial
multiplicities, beginning with periodic or quasi-polynomial weights and then moving
to base zeta functions with several or higher-order poles.

That is the precise target here. The extension is not obtained by assuming that an
arbitrary arithmetic Dirichlet series continues meromorphically. Instead, we state
such a continuation as a hypothesis in the general theorem, prove the requisite
normal convergence, and verify the hypotheses unconditionally in explicit families.
For these families the result is a complete coefficient algorithm rather than an
asymptotic extrapolation.

The new scope relative to \cite{prior} consists of a general finite-principal-part
transfer theorem, the separation of the mean and oscillatory parts of
quasi-polynomial multiplicities, independent Hurwitz factors with all their
Stieltjes residue layers, explicit Gamma reciprocity and primitive formulas, and
an all-order finite compensation for removing a polylogarithmic twist. The
alternating identities extend the report's simple-pole alternants to actual
arithmetic higher poles. The argument is elementary once the correct germ is
retained; its simplicity is not a reason to identify a Laurent coefficient with
an undefined derivative.

\subsection{Classical inputs and the limits of a novelty assertion}
The Hurwitz continuation, its parameter derivative, and its relation to
$\log\Gamma$ are standard \cite{dlmf-hurwitz}. The Pochhammer expansion, the
Gamma product, and the polylogarithm expansion near unity are also classical
\cite{dlmf-gamma,dlmf-polylog}. First-derivative multiplicative-anomaly formulas
have an extensive prior theory; Dowker's treatment includes products of linear
factors and polynomial degeneracies \cite{dowker}. We use that theory as context,
not as a new result. The present proofs establish the stated extensions relative
to the inspected repository. The literature inspection does not prove that every
individual consequence is absent from all earlier publications.

The source audit is targeted. The canonical chapters, the relevant earlier
spectral-product section and its research questions, and the incoming-directory
inventory were inspected. The incoming ZIP files were not all extracted and
audited. The artifact records the source revision and the precise question being
answered; it does not assert a line-by-line audit of the entire repository.

\subsection{Laurent jets and local coefficients}
Throughout, $[s^r]F(s)$ denotes a Laurent coefficient. When $r\ge0$, define
\begin{equation}\label{asj:eq:jetdef}
 \calJ_r F=r![s^r]F(s).
\end{equation}
This equals $F^{(r)}(0)$ only when $F$ is holomorphic there. A ``first jet''
means $[s]F(s)$ under this convention, not the finite part of $F'(s)$ under a
separately chosen change of coordinate. The regulator coordinate $s$ is fixed.

For an expansion at infinity, $[x^{-p}]f(x)$ means the coefficient in the
Laurent series in $x^{-1}$. In particular, $[x^{-1}]$ is the negative of the
complex-analytic residue at infinity when that terminology applies. This
coefficient is computed from a finite truncation in all formulas below.

We use the standard Stieltjes normalization
\begin{equation}\label{asj:eq:stieltjes}
 \zeta(1+s,b)=\frac1s+\sum_{n\ge0}\frac{(-1)^n}{n!}\gamma_n(b)s^n,
 \qquad b>0.
\end{equation}
Thus $\gamma_0(b)=-\psi(b)$ and $\gamma_0(1)=\gamma$, Euler's constant.
Spectral superscripts on $\zeta$ differentiate its first argument; the
superscript on $\psi^{(k)}$ is an argument derivative. All logarithms of positive
real quantities are real. The proofs also give local holomorphic shift
continuations on any domain where compatible logarithms have been fixed.

\subsection{What is not being claimed}
An exact polynomial identity in the exponent variables is not an independence
theorem for its special-function coefficients. A finite numerical residual is
not a proof of equality. Neither a new Gamma-series coordinate nor a zeta-jet
coordinate is claimed to be irreducible in a smaller arithmetic algebra.
The ordinary proofs here have not been formalized in Lean or another proof
assistant.
